\begin{result}{F.3.26}{The universal quotient on full symplectic basis tangencies in every genus}{Proved; exact support and normalizer maps}
Let $g\geq3$, let $V$ have symplectic basis
$x_1,y_1,\ldots,x_g,y_g$ over $\Q_\ell$, and put
$z_i=[x_i,y_i]$, $\omega=\sum_i z_i$.
Write $\mathfrak h_{g,6}$ for the special degree-six
derivations of $\mathbb L(V)$, and $U_{g,6}$ for the
universal arithmetic image with its canonical special
representative and the actual quotient
\[
                 q_{g,6}:U_{g,6}\longrightarrow Q_{g,6}
\]
of Nakamura, as in F.3.24~\cite[pp.~115--117]{AuditNakamura1996}. Its character is the cube
of the similitude multiplier. Define $T_{g,\mathrm{all}}$
by tangency to every one of the specified basis lines:
\[
 D(x_i)=[A_i,x_i],\quad D(y_i)=[B_i,y_i],
 \qquad A_i,B_i\in\mathbb L_6(V).
\]
Then
\[
                 q_{g,6}(U_{g,6}\cap T_{g,\mathrm{all}})=0.
                                                               \tag{AG.1}
\]
This statement retains the full homology space of the
specified curve. A proper CM summand supplies only its
own line constraints and is not substituted for $V$.

\end{result}
\begin{proof}
We reduce every balanced tangency to its exact support on at most four handles and retain the actual universal submodule throughout.
\paragraph{1. The actual geometric-normalizer equation in every genus.}
The construction of F.3.24 uses the moduli of a genus-$g$
curve with one point and one nonzero tangent vector.
Its boundary-fixed geometric monodromy has proalgebraic
image $R_g$ in automorphisms of the rational punctured
Malcev group. The actual profinite geometric image lies
in $R_g(\Q_\ell)$ at each finite nilpotent level by
polynomial continuity and the dense discrete image.
Its algebraic graded image is the Johnson image $J_m$:
Hain's original Theorem~3.4, p.~605, explicitly applies
for all $g\geq3$; Theorem~4.10, p.~611, identifies
the Torelli Malcev weight filtration with its lower
central series for $g\ne2$; Remark~4.11, p.~613,
and the boundary action in Section~14.4, p.~643,
give the strict action maps. Together with the exact
degree-one generation in Enomoto--Satoh, Theorem~3.5
and Proposition~3.6, p.~640, these are the same
typed maps used in F.3.24, now in genus~$g$~\cite{AuditHain1997Torelli,AuditEnomotoSatoh2014}.

Tangential arithmetic monodromy normalizes this actual
algebraic image. Removing the lower inner boundary
term by $\exp(-c\operatorname{ad}Z)$ remains inside
$R_g$, because the boundary Dehn twist is a geometric
element and its powers are Zariski dense in that
one-parameter group. Here $Z$ is the full boundary
logarithm with initial term $\omega$. The free-Lie
centralizer removes the remaining lower inner terms.
Thus the canonical special representative satisfies
\[
       [D,J_1]\subset J_7\qquad(D\in U_{g,6}),
       \qquad J_1=\Lambda^3V.                              \tag{AG.2}
\]
These are leading terms of the actual normalizer
action at a finite nilpotent level. No closure of
an individual discrete filtration kernel is inferred.

Use the exact contraction and cyclic-coinvariant
map $\operatorname{ES}_7$ of F.3.24. It vanishes on
$J_7$ for every $g\geq3$. For $g\geq9$ this is
the stated stable scope of Enomoto--Satoh. For
$3\leq g<9$, the actual boundary-fixed inclusion
of surfaces from genus~$g$ to genus~9 extends
derivations by zero. Both the Johnson map and
the contraction commute with this inclusion;
the cyclic-tensor projection killing the added
letters is a left inverse. Consequently
\[
     \operatorname{ES}_7([D,\delta])=0
      \quad(D\in U_{g,6},\ \delta\in\Lambda^3V).           \tag{AG.3}
\]

\paragraph{2. The complete balanced tangent space has support at most four.}
Give the full diagonal torus its characters
$x_i:a_i$, $y_i:t/a_i$. The weight-$t^3$ subspace
$T_g$ of $T_{g,\mathrm{all}}$ has unique conjugators
$A_i,B_i$ of degree~6 and weight $t^3$. They are
balanced: every handle occurs equally often through
its two generators, and the sum of the handle
multiplicities is three. The equation $D\omega=0$
therefore splits into blocks indexed by
\[
               n=(n_1,\ldots,n_g),\qquad \sum_i n_i=4.
\]
Every nonzero $n_i$ is positive, so each block involves
at most four handles. Up to permutation its possible
types, sizes, and rational ranks are
\[
\begin{array}{c|r|r|r}
\text{positive entries of }n&\text{rows}&\text{columns}&\text{rank}\\\hline
(4)&8&6&6\\
(3,1)&140&66&66\\
(2,2)&312&120&119\\
(2,1,1)&1260&360&359\\
(1,1,1,1)&5040&960&958.
\end{array}                                               \tag{AG.4}
\]
The first four blocks are the exact matrices already
certified in F.3.22 and F.3.24. The last is the only
new block, certified by
\texttt{code/check\_genus\_four\_tangent\_block.py}.
Its ordered $958$-minor has determinant~1 modulo~101.
Two exact independent integer kernels give the
matching upper bound; their independence minor is
$-1$. All matrix entries are coefficients of standard
Lyndon commutators in the free associative algebra.
Thus this finite calculation proves the full rational
rank, rather than only a numerical approximation.

Here is the resulting complete basis description.
There is one $P_{ij}$ on each pair of handles, with
common conjugators on each paired pair of generators
\[
 A_i=B_i=[z_j,[z_i,z_j]],\quad
 A_j=B_j=-[z_i,[z_i,z_j]],
\]
and zero conjugators elsewhere. On each triple
$\{i,j,k\}$ there are three $T_{i;jk}$, indexed
by the repeated handle $i$, with
\[
 \begin{split}
 A_i=B_i&=[z_j,[z_i,z_k]]+[z_k,[z_i,z_j]],\\
 A_j=B_j&=-[z_i,[z_i,z_k]],\\
 A_k=B_k&=-[z_i,[z_i,z_j]].
 \end{split}                                               \tag{AG.5}
\]
On each four-element set there are two elements
$Q^{(1)},Q^{(2)}$. Number those four handles
$1,2,3,4$, and let $L_{ijk}$ be the standard
Lyndon bracket for the word $ijk$ in formal
ordered letters $z_1,z_2,z_3,z_4$. Its values
for the orders used below are fixed by
$L_{ijk}=[z_i,[z_j,z_k]]$ for $i<j<k$ and
$L_{ikj}=[[z_i,z_k],z_j]$ for $i<j<k$.
The common conjugators of the two generators
on each handle are the following columns:
\[
\begin{array}{c|c|c}
 &Q^{(1)}&Q^{(2)}\\\hline
1&L_{234}+L_{243}&-L_{243}\\
2&L_{143}&-L_{134}-L_{143}\\
3&-L_{142}&-L_{124}\\
4&L_{123}&L_{132}.
\end{array}                                               \tag{AG.6}
\]
The program verifies $\sum_i[A_i,z_i]=0$ as an
exact associative-word identity for each column,
then substitutes $z_i=[x_i,y_i]$ and verifies the
two corresponding full $960$-coordinate kernel
vectors. Since the displayed elements occupy
different blocks except for the independent pair
on each quadruple, (AG.4) proves that all these
elements form a basis of $T_g$. Its dimension is
\[
                  {g\choose2}+3{g\choose3}+2{g\choose4}.
                                                               \tag{AG.7}
\]

\paragraph{3. Exact restriction of the trace to three handles.}
Fix a three-handle set $S$, let $W$ be its symplectic
span, and let $p:\mathbb L(V)\to\mathbb L(W)$ kill
the other generators. Put $D_S=pD|_W$.
It is special, since $p(D\omega)=D_S(\omega_W)$;
it remains tangent to the six basis lines of $W$.
For $\delta\in\Lambda^3W$, extend the corresponding
degree-one derivation by zero outside $W$.
On words in $W$ the projected bracket is precisely
$[D_S,\delta]$.

We check the extra contraction terms from every
outside handle $(x_i,y_i)$ explicitly. Because
$\delta(x_i)=\delta(y_i)=0$, one has
\[
 [D,\delta](x_i)=-[\delta A_i,x_i],\qquad
 [D,\delta](y_i)=-[\delta B_i,y_i].
\]
In the dual tensor
$x_i\otimes[D,\delta](y_i)
 -y_i\otimes[D,\delta](x_i)$, contraction of
the first two slots followed by projection to
$W$ gives
\[
                           p\delta A_i+p\delta B_i.
\]
Indeed the surviving terms are $y_i\delta B_i$
and $x_i\delta A_i$, with the signs of the two
symplectic pairings cancelling the displayed
dual-tensor signs. Every other term retains an
outside letter and projects to zero. Each
displayed term is a Lie polynomial of degree~7
in $W$. It is zero in cyclic coinvariants,
because every Lie bracket is an associative
commutator and its two products are cyclic
rotations. Therefore the exact identity is
\[
 p_{\mathrm{cyc}}\operatorname{ES}_{7,V}([D,\delta])
       =\operatorname{ES}_{7,W}([D_S,\delta]).              \tag{AG.8}
\]
This calculation needs only tangency; equality
$A_i=B_i$ was not assumed in proving it.

For $D\in U_{g,6}\cap T_g$, (AG.3) and (AG.8)
put $D_S$ in the exact genus-three trace kernel
computed in F.3.24. Every basis element in
(AG.5)--(AG.6) whose support is not contained
in $S$ restricts to zero: for an inside active
generator its conjugator contains a letter
from every other handle in its support, while
an outside active generator is simply omitted.
Hence, writing the complete expansion of $D$
in the basis above, F.3.24 gives
\[
    \text{every coefficient of }P_{ij}=0,
    \qquad
    c_{i;jk}+c_{j;ik}+c_{k;ij}=0
    \text{ on each triple}.                               \tag{AG.9}
\]
Each pair belongs to a three-handle set because
$g\geq3$. Thus the remaining triple contributions
are differences of the three $T$ elements.

\paragraph{4. The actual quotient kills the remaining components.}
Every handle permutation is an actual element of
$\operatorname{Sp}(V)$. It permutes the three
$T$ elements on a fixed triple transitively.
Therefore every symplectic-invariant linear
functional has the same value on those three
elements and vanishes on their differences.

On a fixed quadruple the program verifies the
stronger exact identities
\[
       \sum_{\pi\in\mathfrak S_4}\pi Q^{(1)}=0,
       \qquad
       \sum_{\pi\in\mathfrak S_4}\pi Q^{(2)}=0.            \tag{AG.10}
\]
The action permutes both the handle indices of
the conjugators and every formal $z_i$ inside
them. These are integer associative-word
identities, checked before the substitution
$z_i=[x_i,y_i]$. There is also a direct independent
proof of these averages. Fix an output handle. For
any input conjugator $A_i$, summing the permutations
that carry handle $i$ to that output sums all six
permutations of its multilinear degree-three Lie
word on the remaining three $z$ variables. Write
that word as a sum of nested brackets. Each nested
bracket has zero such sum, by pairing permutations
that exchange the two arguments of its inner
bracket. Sum this identity over the four possible
input handles. It gives (AG.10) without a rank or
representation-theoretic inference. If a functional
is invariant, its value on the left side is $24$
times its value on the original vector. The coefficient
field has characteristic zero, including when
$\ell=2$ or~3; hence the functional vanishes
on both quadruple basis elements.

The actual submodule $U_{g,6}$ is
$\operatorname{Sp}(V)$-stable, and $q_{g,6}$
is symplectic-invariant. Complete reducibility
of finite-dimensional symplectic representations
extends it to an invariant functional on
$\mathfrak h_{g,6}$. This is only a linear
extension; it assigns no arithmetic meaning
to a vector outside $U_{g,6}$. Equations
(AG.9)--(AG.10) show that the extension
vanishes on $D$, proving
\[
                         q_{g,6}(U_{g,6}\cap T_g)=0.
                                                               \tag{AG.11}
\]

Finally, both $U_{g,6}$ and $T_{g,\mathrm{all}}$
are stable under the full diagonal similitude
torus. For $U_{g,6}$, symplectic stability and
the scalar action $s^6$ on degree-six derivations
give similitude stability through the central
isogeny. The explicit Lagrange weight projector
of F.3.25 onto the $t^3$ eigenspace preserves
their intersection and satisfies $q_{g,6}e=q_{g,6}$.
Applying (AG.11) proves (AG.1). This projection
is applied after an actual arithmetic class
has satisfied all line tangencies; no lift
of this subsequent linear operator is needed.

\paragraph{Remaining proper-summand map.}
The proof has established a linear restriction
map $D\mapsto D_S$ and the precise trace identity
(AG.8). It has not identified $q_{g,6}(D)$ with
the degree-six arithmetic quotient of a proper
symplectic summand. Such an identification, or
sufficient additional mixed Johnson equations,
is still needed when the actual conjugate CM
normalizer lines fill only part of the homology
of the detecting curve in F.3.23.

\end{proof}
